\Askhsh{34324}{4}
\begin{ylh}{1}
    \item 3.11. Ανισοτικές σχέσεις πλευρών και γωνιών
    \item 3.14. Σχετικές θέσεις ευθείας και κύκλου
    \item 3.15. Εφαπτόμενα τμήματα
    \item 4.2. Τέμνουσα δύο ευθειών - Ευκλείδειο αίτημα
    \item 5.9. Μια ιδιότητα του ορθογώνιου τριγώνου
    \item 5.10. Τραπέζιο
    \item 5.11. Ισοσκελές τραπέζιο.
\end{ylh}
% ===================================================================
%  Αρχείο: 34324.tex (Fragment)
%  Αυτόματη μετατροπή μέσω doc2latex.py
% ===================================================================

ΘΕΜΑ 4

Δίνεται κύκλος κέντρου Ο και ακτίνας ρ. Από σημείο Α εξωτερικό του κύκλου θεωρούμε τις εφαπτόμενες του κύκλου που εφάπτονται σε αυτόν στα σημεία Β, Γ και τέτοιες ώστε, η γωνία $Β\widehat{Α}\Gamma$ που σχηματίζουν τα εφαπτόμενα τμήματα ΑΒ και ΑΓ να είναι $60^{ο}$. Έστω ότι η ευθεία ΑΟ τέμνει τον κύκλο στο σημείο Δ και η εφαπτόμενη του κύκλου στο Δ τέμνει τα τμήματα ΑΒ και ΑΓ στα σημεία Ζ και Ε αντίστοιχα.

Να αποδείξετε ότι:

\begin{alist}
\item ΟΑ=2ρ, \monades{6}

\item το τρίγωνο ΑΖΕ είναι ισόπλευρο, \monades{5}

\item ΑΖ = 2ΖΒ, \monades{7}

\item το τετράπλευρο ΕΖΒΓ είναι ισοσκελές τραπέζιο. \monades{7}


\begin{center}
\begin{tikzpicture}[scale=1]
\tkzDefPoint(0,0){O}
\tkzDefPoint(4,0){A}
\tkzDefPoint(2,0){D} % Point D is on OA, radius is 2 since OA=2R and angle is 30.
\tkzDrawCircle(O,D)

\tkzDefTangent[from=A](O,D) \tkzGetPoints{B}{C}

% Tangent at D
\tkzDefLine[orthogonal=through D](O,A) \tkzGetPoint{tangent_D_dir}
\tkzInterLL(D,tangent_D_dir)(A,B) \tkzGetPoint{Z}
\tkzInterLL(D,tangent_D_dir)(A,C) \tkzGetPoint{E}

\draw[l3] (A)--(O);
\draw[l3] (A)--(B);
\draw[l3] (A)--(C);
\draw[l3] (Z)--(E);
\draw[l3] (O)--(B);
\draw[l3] (O)--(C);
\draw[l3, dashed] (B)--(C);

\tkzDrawPoints(O,A,B,C,D,Z,E)
\tkzLabelPoint[left](O){$O$}
\tkzLabelPoint[right](A){$A$}
\tkzLabelPoint[above](B){$B$}
\tkzLabelPoint[below](C){$\Gamma$}
\tkzLabelPoint[above right](D){$\Delta$}
\tkzLabelPoint[above right](Z){$Z$}
\tkzLabelPoint[below right](E){$E$}

\tkzMarkRightAngle(O,B,A)
\tkzMarkRightAngle(O,C,A)
\tkzMarkRightAngle(A,D,Z)
\tkzMarkAngle[size=0.6,mark=none](C,A,B)
\tkzLabelAngle[pos=1](C,A,B){$60^\circ$}
\end{tikzpicture}
\end{center}
\end{alist}
